We exemplary applied our cards to one real world use case at our industry partner. However, in order to avoid revealing confidential information we present results of a toy audit here in the paper.
Consider a large bicycle manufacturer producing around 50\,000 pieces per year. The manufacturer plans to upgrade its product portfolio with carbon fiber frames due to their lightness. The number of bicycles with carbon frames is estimated to be around 5\,000 pieces per year. However, carbon frames have one drawback: they can break easily if something went wrong during the production process. For this reason, the manufacturer plans to use a two-fold quality inspection consisting of a regularly reoccurring computed tomography that can detect production errors reliably but at a very high cost every 100 pieces. On top of that, each frame shall undergo a visual quality inspection using computer vision before shipping (binary classification problem into defect and not-defect parts). If the computer vision (CV) tool predicts a defect, a worker will double check the frame.

\newpage
{\small
\begin{longtable}{p{0.15\linewidth}|p{0.35\linewidth}p{0.5\linewidth}}  
    \toprule
    \multicolumn{3}{c}{\bf{Use Case Card Part 1}}\\
    \toprule
    \bf{Step} & \bf{Requirement} & \bf{Measures} \\
    \hline
    \hline
    General 
        & Name contact person                 & John Doe, Project Manager \\\cline{2-3}  
        & List groups of people involved      & John Doe, Jane Roe, Richard Miles \\\cline{2-3}  
        & Summarize the use case shortly      & Produced carbon fiber bicycle frames shall undergo a quality inspection. Using computed tomography (CT) scans yields high precision, but would be associated with high costs. To reduce costs a machine learning model shall perform automated visual quality inspections on every frame and expensive CT-scans will be performed less. If a defect is predicted, the frame will be inspected by a trained professional. To ensure quality standards every \textit{100}-th frame will be inspected with a CT image. \\\cline{2-3}  
        & Describe the status quo             & To this day carbon fiber bicycle frames can only be inspected manually or by CT-scans. Current non carbon fiber frames are stress-tested and inspected during quality end control. During the initial year 5\,000 carbon fiber frames were produced and manually inspected, with 300 defective parts detected. For a larger database, an additional 200 images of defective frames were obtained from a service partner. \\\cline{2-3}  
        & Describe the planed interaction of sub-components (e.g., different software modules or hardware and software)   & After the pressing process fabricated parts will be photographed by industrial high speed cameras and evaluated by the machine learning model. The predictions of the ML model will be visualized with an interface for quality management workers, predicted defects will be highlighted and require a manual inspection result. For each frame produced, a randomized algorithm determines whether a CT-scan should be performed. Shift workers will then initiate the CT-scan process and document the result.   \\\cline{2-3}  
        & Provide a short solution summary    & The business problem can be translated to a binary classification problem, which can be solved by means of industrial CV. \\
        
    \hline
    Problem definition 
        & Clearly describe the learning problem  & Carbon fiber frames are highly vulnerable to tears and deformations, which can result in stability issues and thus, braking of the frame. In use this could cause accidents and harm to cyclists. Manual and CT-scan inspections are too costly to be performed for every frame. A solution that uses algorithms must work accurately, but fast enough not to delay production times. \\\cline{2-3}  
        \pagebreak
        \cline{2-3}  
        & Describe the disadvantages of the current approach & High costs for CT-scans, manual visual inspections are not feasible due to high time consumption and as a mundane task not appealing to workforce.\\\cline{2-3}  
        & Argue why classical approaches are not sufficient & Rule-based algorithms may not yield the necessary confidence. \\
        
    \hline
    Solution approach 
        & Describe the integration into the current workflow           & Images of the frames are automatically taken. The ML model will evaluate the images in a few seconds and present the results to current shift workers via an interface. Detected defective frames will automatically be marked for manual inspection. The pass/fail labels assigned by the model can  be manually overwritten by workforce. Workers will perform manual inspections on defective parts before submitting the frames to further assembly.\\\cline{2-3}  
        & Formulate the learning problem &  Binary label classification with optional heatmaps / attribution maps applied to the predictions.\\\cline{2-3} 
        & Formulate the KPIs for go-live & A recall above 99\% is targeted while preserving a precision of more than 90\% and the maximum inference time should not exceed 1 second.\\\cline{2-3}  
        & Provide a short data description & Photos are taken with an industrial camera. Both sides of the frame are photographed. The pictures are of size 1024 x 1024 pixels.\\
    \hline
\end{longtable}
} 


\newpage
{\small
\begin{longtable}{p{0.15\linewidth}|p{0.35\linewidth}p{0.5\linewidth}}  
    \toprule
    \multicolumn{3}{c}{\bf{Use Case Card Part 2}}\\
    \toprule
    \multicolumn{3}{c}{Risk Assessment} \\
    \hline
    \hline
    Risk Class          
        & Categorize the application into a risk class according to the AI Act & Low risk. Although the carbon fiber frame is a safety relevant part of the bike, the AI is used for quality control and not part of the final product.\\
        
    \hline
    Human agency and oversight 
        & Rate the level of autonomy & Low risk since the AI only gives a suggested quality label. Manual overwrite is always possible. \\
        
    \hline
    Technical robustness and safety 
        & Evaluate the danger to life and health & Miss-classified broken frames might be shipped and customers may be hurt. \\\cline{2-3}  
        & Identify possibilities of non-compliance & No risk, standard safety requirements are fulfilled without the use of an AI, the AI is an additional layer of security. \\\cline{2-3}  
        & Identify customer relevant malfunction & No risk, the AI is solely used for internal quality control. \\\cline{2-3}  
        & Identify internal malfunction & Low risk of production stoppage or higher production costs. \\
        
    \hline
    Privacy and data governance     
        & List risks connected with customer data & No risk (no data collected). \\\cline{2-3}  
        & List risks connected with employee data & Risk of analyzing quality of work of production staff. \\\cline{2-3}  
        & List risks connected with company data & Low risk of production flaws being leaked to the outside world. \\
        
    \hline
    Transparency                    
        & Evaluate effects of incomprehensible decisions or the use of a black box model & Might result in trust issues in the AI within the workforce.                                 This will lead to either more manual inspections and with that higher costs or fewer inspections and thus lower quality.  \\\cline{2-3}  
        
    \hline
    Diversity, non-discrimination and fairness 
        & Check for possible manipulation of groups of people & No risk.  \\\cline{2-3}
        & Check for discrimination of groups of people regarding sensitive attributes  & No risk. \\
        
    \hline
    \pagebreak
    \hline
    Societal and environmental well being 
        & List possible dangers to the environment & No risk. Good quality assurance reduces waste.\\\cline{2-3}
        & Consider ethical aspects & No risk. \\\cline{2-3}  
        & Evaluate effects on corporate actions & An ELSI-Group (Ethical, Legal and Social Implications) has been formed by the workers council. \\\cline{2-3}  
        & Identify impact on the staff & The AI is solely used to ensure quality in the product. No devaluation of workforce quality is practiced. No analyses are conducted on the quality of work of the production staff and no employment contracts are terminated due to the use of AI. \\    
        
    \hline
    Accountability                      
        & Estimate the financial damage on failure & Broken frames might be shipped and need to be replaced. If customers get injured they might file for                             lawsuits. Misclassified flawless frames might get disposed of. \\\cline{2-3}  
        & Estimate the image damage on failure & Shipped misclassified broken frames might associate the company with poor quality. \\
        
    \hline
    Norms               
        & List of relevant norms in this context & \href{https://www.iso.org/ics/43.150/x/}{ISO 43.150 Cycles} \\
    \bottomrule
  \end{longtable}
}



\newpage
{\small
\begin{longtable}{p{0.15\linewidth}|p{0.35\linewidth}p{0.5\linewidth}}  
    \toprule
    \multicolumn{3}{c}{\bf{Data Card}}\\
    \toprule
        
    \bf{Step}            & \bf{Requirement}                                & \bf{Measures} \\
    \hline
    \hline
    
    General        
        & Name originator of the dataset and provide contact person            & John Doe, Project Manager.  \\\cline{2-3}               
        & Describe the intended use of the dataset                     & Train a supervised ML model for the recognition of defective bike frames.   \\\cline{2-3}              
        & Describe licensing and terms of usage                     & Part of the data was recorded internally and can be used without any restrictions. Another part of the data is provided from a service partner under a specific licensing agreement. \\\cline{2-3}             
    
    \hline
    Data-Description  
        & Collect requirements for the data before starting data collection & Before starting the data collection process, domain experts where involved to discuss how defects occur on the bikes. We discussed the proper amount of data and the necessary resolution of the images.\\\cline{2-3}
        & Describe a data point with its interpretation             & One data point consists of two RGB images with a  resolution of 1024$\times$1024 pixels. One image shows the frame frame from the left, the other one from the right. Each image was taken with an industrial camera of type [MODEL].  \\\cline{2-3} 
        & Maybe, provide additional documentation to understand the data & Not necessary, image data are human interpretable.        \\\cline{2-3}
        & If there is GDPR relevant data. describe it  & No personal data.               \\\cline{2-3}
        & If there is biometric data, describe it                   &   No biometric data.               \\\cline{2-3}
        & If there is copyrighted data, summarize it                &   No copyrighted data.               \\\cline{2-3}
        & If there is business relevant information, describe it          & Yes, the data of the intact bike frames provides information about the glueing technique for the carbon fiber frames. The defective bike frames contain information about weaknesses and unstable parts. Both information could be relevant for competitors.        \\
       
    \hline
    Collection        
        & Describe the data collection procedure and the data sources        & Data was systematically collected by our company during the first year as well as by our service partner  using  commercial industrial cameras of type [MODEL] under diffuse lighting conditions for both cases. We collected data by taking the images of the frames directly after production. Our service partner captured images of the frames it got from its customers. We merged the data and visually checked for homogeneity. After that we labeled the data using an external labeling company.  \\\cline{2-3}
        
        \pagebreak
        \cline{2-3}
        & Use data version control                                  & A data versioning system is applied in the cloud.\\\cline{2-3}
        & Consider prior requirements for data                      & The exact amount of data was collected as defined above.\\\cline{2-3}
        & Include and describe metadata                             & We include the type of bicycle frame as well as the exact date the image was taken. \\\cline{2-3}  
        & Involve domain experts and describe their involvement     & Domain experts were involved during the collection of requirements for the dataset. They also provided instructions and a gold standard dataset for the labeling company. They made several spot checks to control for correctness of the labels afterwards.       \\\cline{2-3}
        & Describe technical measures to ensure completeness of data & We performed an embedding followed by a clustering to identify groups of bike frames other than the different bike models. Those clusters were checked by domain experts. The numbers sold of our different bike models are accurately represented in the data. Also we tried to include enough defect data to make sure the trained model learns how a defect looks in the images.    \\\cline{2-3}
        & Describe and record edge cases                            & One edge case is where the use of to much glue could lead to a label defect. This also happened a few times during the labeling process.    \\\cline{2-3}
        & If personal data is used, make sure and document that all individuals know they are part of the data & No individuals are represented in this dataset.        \\\cline{2-3}
        & Describe if and how the data could be misused & It could be used by competitors to find weaknesses and important design choices in our product line.       \\\cline{2-3}
        & If fairness is identified as a risk, list sensitive attributes & No information is saved to reconstruct the identity of workers.  \\\cline{2-3}
        & If applicable, address data poisoning                     &  Not applicable.    \\\cline{2-3}
    
    \hline
    Labeling        
        & If applicable, describe the labeling process              
        &  The labeling company got a short briefing and a gold-standard labeling set by our domain experts. After the labeling we made spot checks on the dataset.        \\\cline{2-3}  
        & If applicable, describe how the label quality is checked  &  Domain experts perform spot checks. \\
            
    \hline
    Splitting        
        & Create and document meaningful splits with stratification                     & We divided the dataset into easy and hard data points. In the hard data points there were new bike models previously not seen by the model.\\\cline{2-3}  
        & Describe how data leakage is prevented                                        & We use the same camera type for all images. After training we perform heatmap analysis to ensure the model focuses on the bike frame.              \\\cline{2-3}  
        & Recommendation: test the splits and variance via cross-validation             & The variance between several evaluations on the test set is below 1\%.    \\\cline{2-3}  
            
        \pagebreak
        \cline{2-3}  
        & Recommendation: split dataset into difficult, trivial and moderate              & An additional (hard) test set was created that contains novel frames from the service provider that are not part of train data. This allows us to evaluate the applicability of our model to future products. \\\cline{2-3}  
        & Recommendation: put special focus on label quality of test data            & The gold-standard data is included in the test set and all test labeles were inspected by two separate experts. \\\cline{2-3}  
        & Reminder: Perform separate data preprocessing on the splits & Done. \\\cline{2-3}  
              
    \hline
    Preprocessing        
        & Document and motivate all processing steps that are a fixed part of the data   &  The image data is normalized from 0 to 1 to ensure a good behaviour of the convolutional neural network (CNN). \\\cline{2-3}  
        & Document whether the raw data can be accessed                                 &  The raw data is stored on our servers and can be accessed at any time by us. It is not publicly available.          \\\cline{2-3}  
        & If sensitive data is available highlight (pseudo)-anonymization               & No (pseudo)-anonymization needed. \\\cline{2-3}  
        & If fairness is a risk, highlight fairness specific preprocessing             & We do not record information to draw conclusion about the identity of specific workers.\\
             
    \hline
    Analyzing        
        & Understand and document characteristics of test and training data     &  We have 4\,700 intact bike frames and 500 bike frames with defects. 300 of the defective frames are from our own production and are identical to our frames, 200 are from the service provider. We use 20$\%$ of the data for the test split.\\\cline{2-3}   
        & Document why the data distribution fits the real conditions or why this is not necessary for the use case         &  Images were taken directly at our test production line and at the location of our service provider using real intact and defective frames.\\\cline{2-3}  
        & Document limitations such as errors, noise, bias or known confounders   &  It is possible that non-superficial defects cannot be identified by a visual check.\\\cline{2-3}  
    \hline
    Serving       
        & Describe how the dataset will be maintained in future &  New images from the production line and the service provider will be added to the dataset after labeling by the external company. \\\cline{2-3}  
        & Describe the storage concept   &   We operate a server with redundancy to store the image data.   \\\cline{2-3} 
        & Describe the backup procedure  &   A compressed and encrypted backup of the current dataset is saved at an external cloud storage provider.      \\\cline{2-3}  
        & If necessary, document measures against data poisoning                &  Not necessary.        \\
                      
    \bottomrule
\end{longtable}
}



\newpage
{\small
  \begin{longtable}{p{0.15\linewidth}|p{0.35\linewidth}p{0.5\linewidth}}
    \toprule
    \multicolumn{3}{c}{\bf{Model Card}}\\
    \toprule
    \bf{Step}   & \bf{Requirement}   &  \bf{Measures} \\
    \hline
    \hline
    General                     
        & Provide Name and details of a contact person 		    & Jane Doe, Project Manager\\\cline{2-3}
        & Provide Name of the person who created the model      & John Doe, ML Engineer\\\cline{2-3}
        & Document the creation date and version of the model   & 01-2023, Version 1.0 \\\cline{2-3}
        & Describe intended use of the model                    & Detection of defects in carbon fiber frames\\
        
    \specialrule{.075em}{0em}{0em}
    Description       
        & Describe the architecture of the used model    	& Custom CNN\\\cline{2-3}
        & Describe the used hyper-parameters         	    & Six convolution layers with ReLu activation followed by max-pooling. Sigmoid activation in the final layer. The model is trained for 200 epochs with learning rate 1e-6.\\\cline{2-3}
        & Document the training, validation, and test error     & Recall on train and both moderate and hard test set is 100\%. However, the precision is only around 90\%. Hence, the recall is good enough to apply the solution but an improved precision would further reduce  manual checks.\\\cline{2-3}
        & Document computation complexity, training time and energy consumption (for foundation models describe steps taken to reduce energy consumption) & Model was trained on cloud resources using 6 NVIDIA A100 GPUs for 32 hours. Estimated energy consumption: 70 kWh.\\
                        
    \specialrule{.075em}{0em}{0em}
    Explainability and interpretability    
        & Document the demand for explainability and interpretability 	& The model should provide heatmaps to workers to help them detect regions of interest.\\\cline{2-3}
        & Describe taken actions if any 			& GradCam (or similar) will be applied post-hoc during operation.\\   

    \specialrule{.075em}{0em}{0em}
    Feature engineering, feature selection and preprocessing   
        & Describe the consideration of explainability and interpretability (if relevant)  & Not relevant because there is no feature engineering.\\\cline{2-3}
        & Describe feature engineering and selection with (domain specific) reasoning   & No feature engineering.\\\cline{2-3}
        & Describe and reason of preprocessing steps  & Data augmentation in form of brightness and Gaussian noise.\\

    \specialrule{.075em}{0em}{0em}
    \pagebreak
    \specialrule{.075em}{0em}{0em}
    Model selection                     
        & Describe the consideration of explanibility, interpretability and fairness (if relevant) 	& A standard heatmap method must be applicable to the model.\\\cline{2-3}
        & Describe the base line model and its evaluation	& Multiple baselines were evaluated, a linear model, ViT and smaller and larger CNNs.\\\cline{2-3}
        & Describe the reasoning for the model choice 		& The final model was chosen based on a trade-off between performance and model size.\\\cline{2-3}
        & Describe why the complexity of model is justified and needed 	& Yes, smaller or simpler models performed worse. Larger models increased the inference time too much.\\\cline{2-3}
        & Document the comparison to other considered models 	  & Using a deeper model such as ResNet only marginally improves performance by less than 1\%.\\\cline{2-3}
        & Document the approach of hyper-parameter optimization      &  Random search. Increases test performance by around 5\%. \\\cline{2-3}
        & Describe the model evaluation                   &  The model is evaluated w.r.t. the recall and precision on both the moderate and hard test and, additionally, tested in production.\\
        
    \specialrule{.075em}{0em}{0em}
    Choice of metrics                   
        & Describe selected metrics and describe reasoning regarding use case and fairness    	& The evaluated metric is recall and an inference time below 0.5 seconds. Fairness is not applicable here.\\\cline{2-3}
        & Define the minimum requirement for use in production (domain specific reasons) 	& Recall above 99\% is the target. Either purely from the model or by including a human in the loop. A precision of 90\% is targeted but not mandatory. \\

    \specialrule{.075em}{0em}{0em}
    Model confidence                   
        & Document and quantify uncertainty of the model    & Model confidence was not assessed, since the model was trained w.r.t recall.\\\cline{2-3}
        & Document approach of dealing with uncertainty     & None, see above.\\

    \specialrule{.075em}{0em}{0em}
    Testing in real world setting       
        & Describe test design 		&  During the initial phase (expected to be the first year) all frames will additionally be checked manually. During this phase, faulty frames will be injected into the line in order to validate the AI solution.\\\cline{2-3}
        & Describe possible risks, edge cases and worst case scenarios and create (or simulate) them if possible & Edge cases are tested both using the different test sets (hard and medium) and through injection of faulty parts.\\\cline{2-3}
        & Describe limitations and shortcomings of the model & Change of production process or architecture of frames might introduce new faults, that might not be detcted by the model.\\\cline{2-3}
        & Describe the test results		& The test is still running but revealed problems of the AI solution under specific lighting conditions\\\cline{2-3}
        & Explain the derived actions 	& The dataset was extended to include such conditions\\
    \bottomrule
  \end{longtable}
}



\newpage
{\small
\renewcommand\arraystretch{1.3}
\begin{longtable}{p{0.15\linewidth}|p{0.35\linewidth}p{0.5\linewidth}}
    \toprule
    \multicolumn{3}{c}{\bf{Operation Card}}\\
    \toprule
    \bf{Step}   & \bf{Requirement}   &  \bf{Measures} \\
    \hline
    \hline
    \multirow{7}{=}{General} 
    Scope and aim of monitoring     
        & Describe monitored components & The CV system as a whole is monitored. Input and output of the model are evaluated separately. Additionally the function of the cameras is monitored.\\\cline{2-3}
        & Assess risks and potential dangers according to Use case card & Undetected faults lead to the shipment of broken frames. False Alarms might lead to ineffective manual inspections and trust issues of personnel.\\\cline{2-3}
        & List safety measures for risks & Monitoring of input and output data. Regularly reoccurring computed tomography enables assessment of model performance. Manual inspections enable the assessment of precision.\\\cline{2-3}
                                    
    \hline
    Operating concept               
        & Create and document utilisation concept & If the CV system reports possible damage of a frame, staff does a standardized thorough visual and tactile examination.\\\cline{2-3}
        & Plan staff training & Each employee who conducts the examination needs to take a training session with test. Staff training of the user interface and AI basics is mandatory. \\\cline{2-3}
        & Determine responsibilities & The shift supervisor is responsible for the examination and dedicated IT  personnel for the CV system.\\\cline{2-3}
        
    \hline
    Autonomy of the application     
        & Document decision-making power of AI  & The AI system checks autonomously and only notifies personnel if faults are detected. The user interface allows checking the system and the images made by the cameras at all times. \\\cline{2-3}
        & Determine process to overrule decisions of the AI & Regularly reoccurring computed tomography tests overrule the CV systems decision. Personnel can overrule the AI system's decision after manual inspection of a frame.\\\cline{2-3}
    
    \hline
    Responsibilities and measures   
        & Component wise: assessed risk, control interval, responsibility, measures for emergency & Risks of system components (CV model, cameras, data processing pipeline and user interface) are assessed separately. Further detailed information is attached to this document.\\\cline{2-3}
    
    \hline
    \pagebreak
    \hline
    Model performance               
        & Monitor input and output & Input: Distribution of pixel values are monitored. Output: Number and distribution of defective frames is monitored over a fixed time period.\\\cline{2-3}
        & Detect drifts in input data & Input data is monitored with torch drifts Kernel MMD drift detector method.\\\cline{2-3}
        & Document metric in use to monitor model performance & Model performance is monitored via precision and recall. If precision is over 95\% or under 85\% personnel will be alerted. If the CV system  misses a defective frame (detected by computed tomography) personnel will be alerted immediately. If the percentage of defective frames increases, or percentage of defective frames decreases over time, personnel is alerted immediately.\\\cline{2-3}
    
    \hline
    AI Interface                    
        & Establish transparent decision making process & Decision making is made transparent with heatmaps and displayed on the user interface. \\\cline{2-3}
        & Establish insight in model performance on different levels & Production staff: current frame and heatmaps displayed. Display of heatmaps of frames identified as defective.
        IT personnel: Camera images and heatmaps analogous to production staff. Additionally, access to logs, performance metrics and monitoring as described.\\\cline{2-3}
    
    \hline
    IT security                     
        & Document individual access to server rooms & Access to server rooms is restricted to authorized personnel only\\\cline{2-3}
        & Set and document needed clearance level for changes to AI/Deployment/Access regulation & Only trained and authorized personnel can make changes to the CV system. Changes need to be verified by a second person to be applied.\\\cline{2-3}
        & Establish and audit ISMS & IT system in which the AI application will be integrated is certified according to ISO 270001. A review of the system will be conducted with an IT-security manager to evaluate the need for re-certification.\\ \cline{2-3}
    
    \hline
    Privacy                         
        & Justify and document use or waiver of a privacy preserving algorithm & Not applicable since no personal data is collected.\\\cline{2-3}
        & If applicable, document privacy algorithm and due changes in the monitoring of output data & Not applicable.\\\cline{2-3}
    
    \hline
    MLOps                           
        & Establish versioned code repository of AI, training and deployment & Version control is implemented with Git and MLflow.\\\cline{2-3}
        & Establish maintenance and update schedule & The ML model is retrained every 3 months (if new data is available) on a blend of old and new data or if a drift is detected.\\\cline{2-3}
        & Set regulation for the retraining of the AI and decision basis for the replacement of a model & The model is updated with new training data. The update goes live, if the recall does not drop and precision increases or remains constant. If performance measures drop twice in a row after retraining, data will be analyzed.\\\cline{2-3}
    
    \hline
    Testing and rollout             
        & Determine responsibilities for updates & Updates are managed by dedicated IT personnel. Changes need to be verified by a second person to be applied.\\\cline{2-3}
        & Document software tests & Software tests are automated and documented.\\\cline{2-3}
        & Set period of time that an update must function stably before it is transferred to live status & New software and model versions run separately on the test environment parallel to the production environment for one week. Performance is evaluated and compared to the old version before deployment.\\
                                    
    \bottomrule
\end{longtable}
}